\title{xLSTM: Extended Long Short-Term Memory}

\begin{abstract}
During the 1990s, the concepts of the constant error carousel and gating mechanisms became foundational to Long Short-Term Memory (LSTM) models. These LSTMs have since maintained their relevance and have significantly contributed to advancements in deep learning, serving as the building blocks for the earliest Large Language Models (LLMs). Despite these achievements, the introduction of Transformers—anchored by scalable, parallelizable self-attention—ushered in a paradigm shift, surpassing LSTMs at scale. In this work, we pose a straightforward question: What are the capabilities of LSTMs in language modeling when they are scaled to billions of parameters, enhanced with contemporary LLM methodologies, while addressing longstanding LSTM limitations? Our approach introduces exponential gating combined with robust normalization and stabilization strategies. Additionally, we reengineer the LSTM memory design to yield: (i) the sLSTM, characterized by a scalar memory and update mechanism alongside an innovative memory blending technique, and (ii) the mLSTM, featuring a matrix memory that supports full parallelization via a covariance-based update rule. These LSTM advancements are implemented within residual block architectures, forming xLSTM blocks that are residually stacked to construct comprehensive xLSTM models. By applying exponential gating and revising memory architectures, xLSTM is capable of achieving competitive or superior results compared to leading Transformers and State Space Models, demonstrating strong performance and scalability.\\
Code available at: \url{https://github.com/NX-AI/xlstm}
\end{abstract}

\section{Introduction}
\label{sec:intro}

The Long Short-Term Memory (LSTM) ideas~\citep{Hochreiter:91,Hochreiter:97nips,Hochreiter:97}, i.e., the constant error carousel and gating, were introduced to overcome the vanishing gradient problem of recurrent neural networks~\citep{Hochreiter:91,Hochreiter:00book}:

\begin{align}
\label{eq:lstm_idea}
  c_t \ = \ f_t \ c_{t-1} \ + \ 
          i_t \  z_t \ , \quad  
          h_t \ = \ o_t \ \psi({c_t}) \ . 
\end{align}

The constant error carousel involves incrementally updating the cell state $c_{t-1}$ (shown in green) through the addition of cell inputs $z_t$, with this process being regulated by sigmoid gating mechanisms (blue). The update is governed by both the input gate $i_t$ and the forget gate $f_t$, whereas the output gate $o_t$ manages the emission from the memory cell, specifically the hidden state $h_t$. To produce the hidden state, the cell state is first passed through a squashing or normalization function $\psi$, and then subjected to output gating.

LSTMs have achieved notable success across a wide range of fields \citep{Hochreiter:01,Hochreiter:07,Schmidhuber:15} and dominated text generation approaches until the introduction of Transformers in 2017~\citep{Vaswani:17}. Their capability has been verified on many sequence-based tasks, including text production~\citep{Graves:13,Karpathy:15blog}, handwriting generation~\citep{Graves:13}, sequence-to-sequence translation~\citep{Sutskever:14nips}, program analysis~\citep{Zaremba:14arxiv}, image captioning~\citep{Karpathy:15,Hossain:19}, code synthesis~\citep{Karpathy:15blog}, rainfall-runoff forecasting~\citep{Kratzert:18,Kratzert:19}, and hydrological flood prediction~\citep{Nearing:24}. Within reinforcement learning, LSTMs are recognized as the top-performing models for sequential tasks, demonstrated in systems such as the AlphaStar agent for StarCraft II~\citep{Vinyals:17short}, the OpenAI Five project for Dota 2~\citep{OpenAI:19}, as well as controllers for plasma in nuclear fusion experiments~\citep{Degrave:22short}. LSTMs distinguish themselves through their proficiency in developing abstractions, efficiently distilling semantic content and retaining it within memory cells~\citep{Karpathy:15blog}. This ability is exemplified by the discovery of number and syntax-selective units~\citep{Lakretz:19}, linguistic-specialized neurons~\citep{Bau:19}, and sentiment-encoding neurons~\citep{Radford:17arxiv}. Despite the emergence of newer models, LSTMs continue to find utility in essential, modern-day tasks~\citep{Degrave:22short,Nearing:24} and have exhibited enduring relevance in the field.

Although LSTMs have achieved remarkable performance, they exhibit three primary shortcomings: (i) They lack the ability to revise previous memory decisions. This is illustrated using the {\it Nearest Neighbor Search} task (further details in Appendix~\ref{sec:appNearestNeighborSearch}). Given a reference vector, the model processes a sequence to identify the most similar vector, yielding its corresponding value at the end. As depicted in the left panel of Figure~\ref{fig:lstmProblems}, LSTM’s mean squared error demonstrates its difficulty in updating the stored value upon encountering a more relevant vector later in the sequence. In contrast, our proposed xLSTM, using exponential gating, mitigates this issue. (ii) Storage constraints, wherein information is compressed into scalar-valued cell states, also present a bottleneck. The {\it Rare Token Prediction} task highlights this challenge. The right panel of Figure~\ref{fig:lstmProblems} presents perplexity outcomes on token prediction for partitions based on frequency from the Wikitext-103 dataset~\citep{Merity:17}. LSTMs suffer in predicting less common tokens due to this restricted storage; our xLSTM alleviates this through a matrix-based memory. (iii) A lack of parallel processing is inherent, owing to memory mixing—specifically, the recurrent dependencies between consecutive time steps enforce sequential execution.

Such limitations have contributed to the popularity of architectures like Transformers~\citep{Vaswani:17} for language modeling. An open question remains: what advancements in language modeling are possible by addressing these shortcomings and scaling LSTMs to match the capacities of today’s Large Language Models?

\section{Extended Long Short-Term Memory}
\label{sec:xLSTM}

In response to the constraints of standard LSTM, the Extended Long Short-Term Memory (xLSTM) introduces two principal changes to the foundational concept from Equation~\eqref{eq:lstm_idea}. These innovations—namely exponential gating and reimagined memory structures—expand the LSTM framework through two newly proposed architectures: (i) the sLSTM (detailed in Section~\ref{sec:sLSTM}), characterized by its use of a scalar memory, scalar update mechanisms, and the concept of memory mixing; and (ii) the mLSTM (outlined in Section~\ref{sec:mLSTM}), which employs a matrix-based memory and updates via a covariance (outer product) operation that facilitates complete parallelization. Exponential gating is a key enhancement in both sLSTM and mLSTM. The mLSTM is designed for parallel execution and, as a result, forgoes memory mixing—eliminating recurrent connections between hidden states. Both models can scale to incorporate multiple memory cells; in the sLSTM, this extension includes mixing memory between cells. The sLSTM may also support multiple heads, maintaining memory mixing solely among cells within each head, and not between different heads, thereby introducing a novel paradigm of memory mixing in conjunction with exponential gating. For mLSTM, expanding to multiple heads or multiple cells yields structurally equivalent configurations.

These LSTM variants can be employed as modular components in residual blocks, referred to as xLSTM blocks (see Section~\ref{sec:xLSTMblock}). Constructing models by stacking these xLSTM blocks in a residual fashion leads to xLSTM-based architectures (see Section~\ref{sec:xLSTMarchitecure}). The overall structure and relationships among components in the xLSTM architecture are illustrated in Figure~\ref{fig:xlstm_sketch}.

\subsection{Review of the Long Short-Term Memory}
\label{sec:orgLSTM}

The foundational LSTM model~\citep{Hochreiter:91,Hochreiter:97nips,Hochreiter:97} proposed the scalar memory cell as a core unit responsible for both computation and storage, which mitigates the issue of vanishing gradients~\citep{Hochreiter:91,Hochreiter:00book} by maintaining the error signal through the mechanism known as the constant error carousel (cell state updates). Within this memory cell, three gates are employed—namely, the input gate, output gate, and forget gate. The forget gate was subsequently incorporated by \citet{Gers:00}. The update equations for the LSTM memory cell at time $t$ are as follows:

\begin{align}
c_t \ &= \  f_t \ c_{t-1} \ + \ i_t \ z_t &  & &\text{cell state} \\
h_t  \ &= \ o_t \ \tilde{h}_t \ , & \tilde{h}_t \ &= \ \psi \left( c_t\right)
  &\text{hidden state} \\
z_t \ &= \ \varphi \left( \tilde{z}_t \right) \ , 
  &\tilde{z}_t \ &=  \ \mathbf{w}^\top_{z} \ \mathbf{x}_t \ + \
  r_{z}  h_{t-1} \ + \  b_{z} \ \
  &\text{cell input} \\
i_t \ &= \ \sigma \left( \tilde{i}_t  \right) \ , 
  &\tilde{i}_t \ &= \ \mathbf{w}^\top_{i} \ \mathbf{x}_t \ + \
  r_{i}  \ h_{t-1} \ + \  b_{i} \ \
  &\text{input gate} \\
f_t \ &= \ \sigma \left(  \tilde{f}_t \right) \ , 
  &\tilde{f}_t \ &= \ \mathbf{w}^\top_{f} \ \mathbf{x}_t  \ + \
  r_{f}  \ h_{t-1} \ + \  b_{f} \ \
  &\text{forget gate} \\
o_t \ &= \ \sigma \left( \tilde{o}_t \right) \ , 
  &\tilde{o}_t  \ &= \ \mathbf{w}^\top_{o} \ \mathbf{x}_t \ + \
  r_{o}  \ h_{t-1} \ + \  b_{o} \ \
  &\text{output gate} 
\end{align}

The weight vectors $\mathbf{w}_{z}$, $\mathbf{w}_{i}$, $\mathbf{w}_{f}$, and $\mathbf{w}_{o}$ represent the input connections from $\mathbf{x}_t$ to the cell input, input gate, forget gate, and output gate, respectively. Correspondingly, $r_{z}$, $r_{i}$, $r_{f}$, and $r_{o}$ denote the recurrent weights linking the hidden state $h_{t-1}$ to the cell input and each gate (input, forget, and output). The parameters $b_{z}$, $b_{i}$, $b_{f}$, and $b_{o}$ are the associated biases for these pathways. Activation functions $\varphi$ and $\psi$ are applied to the cell input and the hidden state; both functions are generally taken to be $\tanh$. The $\psi$ function serves to constrain the range of the cell state, preventing it from growing without bound. All gates utilize the sigmoid function for activation, specifically $\sigma(x) = 1/(1+\exp(-x))$. In subsequent versions, the model incorporates a vector of scalar memory cells, $\mathbf{c}_t \in \mathbb{R}^d$, which allows the recurrent transformation via weight matrices $\mathbf{R} \in \mathbb{R}^{d \times d}$ that combine the outputs across multiple memory cells~\citep{Greff:15}; more detailed discussion is provided in Appendix~\ref{sec:apporgLSTM}. Ablation experiments have demonstrated that each part of the memory cell plays an essential role in its functioning~\citep{Greff:15}.

\subsection{sLSTM}
\label{sec:sLSTM}

To enhance LSTMs with a mechanism for revisiting and updating storage choices, we incorporate exponential gates (shown in red), complemented by normalization and stabilization techniques. Specifically, we employ exponential activation functions for both input and forget gates. For normalization purposes, we define a normalizer state which accumulates the products of the input gate and all subsequent forget gates.

The scalar sLSTM forward pass is:

\begin{align}
\label{eq:slstmforward}
c_t \ &= \  f_t \ c_{t-1} \ + \ i_t \ z_t & & 
 &\text{cell state} \\
n_t \ &= \  f_t \ n_{t-1} \ + \ i_t  & & 
 &\text{normalizer state} \\
h_t \ &= \ o_t \ \tilde{h}_t \ ,
  &\tilde{h}_t \ &= \ c_t / n_t
&\text{hidden state} \\
z_t \ &= \ \varphi \left( \tilde{z}_t \right) \ , 
  &\tilde{z}_t \ &=  \ \mathbf{w}^\top_{z} \ \mathbf{x}_t \ + \
  r_{z}  h_{t-1} \ + \  b_{z}
  &\text{cell input} \\
i_t \ &= \ \!\exp \left( \tilde{i}_t  \right) \ , 
  &\tilde{i}_t \ &= \ \mathbf{w}^\top_{i} \ \mathbf{x}_t \ + \
  r_{i}  \ h_{t-1} \ + \  b_{i}
  &\text{input gate} \\
f_t \ &= \ \sigma \left(  \tilde{f}_t \right) \ \text{OR} \
\!\exp \left(  \tilde{f}_t \right) \ ,
  &\tilde{f}_t \ &= \ \mathbf{w}^\top_{f} \ \mathbf{x}_t  \ + \
  r_{f}  \ h_{t-1} \ + \  b_{f}
  &\text{forget gate} \\
o_t \ &= \ \sigma \left( \tilde{o}_t \right) \ , 
  &\tilde{o}_t  \ &= \ \mathbf{w}^\top_{o} \ \mathbf{x}_t \ + \
  r_{o}  \ h_{t-1} \ + \  b_{o}
  &\text{output gate} 
\end{align}

We transfer the original LSTM gating techniques, i.e., input- and/or hidden-dependent gating plus bias term, to the new architectures. Exponential activation functions can lead to large values that cause overflows. Therefore, we stabilize gates with an additional state \mbox{$m_t$~\citep{Milakov:18arxiv}}:

\begin{align}
\label{eq:slstmstabil}
m_t \ &= \ \max \left( \log ( f_t ) + m_{t-1} , \log ( i_t ) \right) &\text{stabilizer state} \\
i'_t \ &= \ \!\exp \left( \log \left ( i_t \right) - m_t \right) \ = \ \!\exp \left( \tilde{i}_t - m_t \right) \ \, 
  &\text{stabil. input gate} \\
f'_t \ &= \ \!\exp \left( \log \left( f_t \right) + m_{t-1} - m_t \right) \ \, 
  &\text{stabil. forget gate}
\end{align}

In Appendix~\ref{sec:appsLSTM}, we demonstrate that substituting $f_t$ with $f'_t$ and $i_t$ with $i'_t$ during the forward computation does not impact the overall network output or alter the gradients of the loss with respect to the model parameters.

\paragraph{New Memory Mixing.} Similar to the original LSTM architecture, sLSTM supports several memory cells (refer to Appendix~\ref{sec:appsLSTM}). The use of multiple memory cells facilitates memory mixing through recurrent matrices $\mathbf{R}_{\mathbf{z}}$, $\mathbf{R}_{\mathbf{i}}$, $\mathbf{R}_{\mathbf{f}}$, and $\mathbf{R}_{\mathbf{o}}$, which create connections from the hidden state $\mathbf{h}$ to the inputs of memory cell $\mathbf{z}$ as well as to the gate vectors $\mathbf{i}$, $\mathbf{f}$, and $\mathbf{o}$. A distinct characteristic in this mixing mechanism is the incorporation of exponential gating. In sLSTM, it is possible to implement multiple heads, each allowing memory mixing internally but preventing interaction across different heads. This combination of head structures and exponential gating brings forth an innovative approach to mixing memory in sLSTM networks.

\subsection{mLSTM}
\label{sec:mLSTM}

In order to improve the memory capacity of LSTMs, we replace the conventional scalar memory cell $c \in \mathbb{R}$ with a matrix memory cell $\mathbf{C} \in \mathbb{R}^{d \times d}$. This allows retrieval operations to be executed as matrix multiplications. Specifically, at each time step $t$, we aim to store a key-value pair consisting of two vectors: the key $\mathbf{k}_t \in \mathbb{R}^d$ and the value $\mathbf{v}_t \in \mathbb{R}^d$, adopting terminology from the Transformer architecture. Subsequently, at a later time $t + \tau$, the stored value $\mathbf{v}_t$ is accessed using a query vector $\mathbf{q}_{t+\tau} \in \mathbb{R}^d$. This configuration follows the paradigm of Bidirectional Associative Memories (BAMs) \citep{Kohonen:72,Anderson:72,Nakano:72,Anderson:77}.

The covariance update rule~\citep{Sejnowski:77,Dayan:91} for storing a key-value pair is

\begin{align}
\mathbf{C}_t \ &= \ \mathbf{C}_{t-1} \ + \ \mathbf{v}_{t} \ \mathbf{k}_t^\top \ .
\end{align}

We consider inputs to be normalized via layer-norm prior to their projection into keys and values, ensuring these vectors possess zero mean. The process of updating the covariance adheres to the optimal rule~\citep{Dayan:91} for maximizing the distinguishability of retrieved binary representations, or equivalently, optimizing the signal-to-noise ratio. It is noteworthy that if retrieval is restricted to pairwise associations and the quadratic computational expense characteristic of attention mechanisms is permitted~\citep{Krotov:16,Krotov:17,Ramsauer:21}, even greater separability can be attained. Notably, this covariance-based approach aligns with the Fast Weight Programmers paradigm \citep{Schmidhuber:92ncfastweights,Schlag:21}, which in subsequent developments included the use of a fixed decay factor applied to $\mathbf{C}_{t-1}$ and a fixed learning rate for scaling 
$\mathbf{v}_{t} \mathbf{k}_t ^\top$~\citep{Ba:16}.
Following this methodology, we incorporate the covariance update procedure into the LSTM architecture; specifically, the forget gate implements the decay operation, the input gate determines the learning rate, and the output gate modulates the magnitude of the retrieved vector.

For this matrix memory, the normalizer state is the weighted sum of key vectors, where each key vector is weighted by the input gate and all future forget gates. Again, the normalizer state keeps record of the strength of the gates. Since the dot product between query and normalizer state can be close to zero, we use the absolute value of this dot product and lower bound it by a threshold (typically 1.0) as done previously~\citep{Sun:23arxiv}. The mLSTM forward pass is:

\begin{align}
\label{eq:mlstm_recurrent_begin}
\mathbf{C}_t \ &= \  f_t \ \mathbf{C}_{t-1} \ + \ 
  i_t \ \mathbf{v}_t \ \mathbf{k}_t^\top & &  &\text{cell state} \\
\mathbf{n}_t \ &= \  f_t \ \mathbf{n}_{t-1} \ + \ 
  i_t \ \mathbf{k}_t & &  &\text{normalizer state} \\
\mathbf{h}_t  \ &= \ \mathbf{o}_t \ \odot \ \tilde{\mathbf{h}}_t
\ , \qquad \qquad 
\tilde{\mathbf{h}}_t \ = \   \mathbf{C}_t \mathbf{q}_t \ / \ 
  \max \left\{ |\mathbf{n}_t^\top \mathbf{q}_t|, 1 \right\} 
   &&&\text{hidden state} \\
\mathbf{q}_t \ &= \ \mathbf{W}_q \ \mathbf{x}_t \ + \ \mathbf{b}_q  & &  &\text{query input} \\
\mathbf{k}_t \ &= \ \frac{1}{\sqrt{d}} \mathbf{W}_k \ \mathbf{x}_t \ + \ \mathbf{b}_k  & &  &\text{key input} \\
\mathbf{v}_t \ &= \ \mathbf{W}_v \ \mathbf{x}_t \ + \ \mathbf{b}_v  & &  &\text{value input} \\
i_t \ &= \ \!\exp \!  \! \left( \tilde{i}_t  \right) \ , 
 \qquad \qquad \ \ \ \ \,
  \tilde{i}_t \ = \ \mathbf{w}^\top_{i} \ \mathbf{x}_t \ + \  b_{i} \ \ \ 
  &&&\text{input gate} \\
f_t \ &= \ \sigma \!  \left(  \tilde{f}_t \right) \ \text{OR} \ \!\exp \!  \! \left(  \tilde{f}_t \right) , \ \ \: \!
\tilde{f}_t \ = \ \mathbf{w}^\top_{f} \ \mathbf{x}_t  \ + \
  b_{f}
  &&& \text{forget gate} \\
\label{eq:mlstm_recurrent_end}
\mathbf{o}_t \ &= \ \sigma \left( \tilde{\mathbf{o}}_t \right) \ , \qquad \qquad \quad \ \ \ \,
  \tilde{\mathbf{o}}_t  \ = \ \mathbf{W}_{\mathbf{o}} \ \mathbf{x}_t \ + \
  \mathbf{b}_{\mathbf{o}}
  &&&\text{output gate} 
\end{align}

Like traditional LSTM models, mLSTM is capable of utilizing several memory cells. In the context of mLSTM, the use of multiple heads serves the same function as multiple cells, since there is an absence of memory interaction or mixing. To address the instability that may arise from the exponential gating in mLSTM, we adopt the stabilization methods outlined for sLSTM (see Equation~\ref{eq:slstmstabil}). Given that mLSTM does not incorporate memory mixing, its recurrence can be expressed in a parallelized manner. Additional information can be found in Appendix~\ref{sec:appmLSTM}.

\subsection{xLSTM Architecture}
\label{sec:xLSTMarch}

\paragraph{xLSTM Blocks.}
\label{sec:xLSTMblock}
The primary function of an xLSTM block is to produce a nonlinear, high-dimensional representation of preceding sequence elements, thereby more effectively differentiating between diverse histories or contexts. Achieving this separation is essential for reliably forecasting the subsequent element in a sequence, such as the next token. To motivate our approach, we appeal to Cover’s Theorem~\citep{Cover:65}, which asserts that nonlinear embeddings into higher-dimensional spaces enhance the likelihood of linear separability among patterns compared to those in the original input space. Our exploration focuses on two types of residual block designs: (i) A residual block utilizing post up-projection (in the spirit of Transformers), where information from the past is first nonlinearly condensed within the original dimension, then elevated via a linear projection into a higher-dimensional space. After a nonlinear activation is applied, a linear map returns the information to its initial dimension; refer to the left side of Figure~\ref{fig:Backbones} and the third column in Figure~\ref{fig:xlstm_sketch}. A more granular illustration is available in Figure~\ref{fig:appsLSTM_detailed} of the appendix. (ii) A residual block employing pre up-projection (similarly to State Space Models), in which the input is first projected linearly into a higher-dimensional space,

encodes the historical information in a non-linear manner within the high-dimensional space before projecting it linearly back to the original representation. When an xLSTM block integrates an sLSTM, we apply the post up-projection component. In contrast, for xLSTM blocks utilizing an mLSTM, the pre up-projection is employed due to the increased memory capacity available in the high-dimensional space. For further clarification, consult the left side of Figure~\ref{fig:Backbones}, the third column of Figure~\ref{fig:xlstm_sketch}, or refer to Figure~\ref{fig:appsLSTM_detailed} in the appendix.

\paragraph{xLSTM Architecture. }
\label{sec:xLSTMarchitecure}
The xLSTM architecture is built by stacking its constituent modules in a residual fashion~\citep{Srivastava:15,He:16}. Our implementation adopts the pre-LayerNorm~\citep{Ba:16b} residual framework, which is the default choice in modern Large Language Models. Further details can be found in the rightmost column of Figure~\ref{fig:xlstm_sketch}.

\subsection{Memory and Speed Considerations}

Unlike Transformers, xLSTM networks maintain linear computational complexity and require constant memory relative to sequence length. The compressive nature of xLSTM memory makes it especially appropriate for use in industrial settings and deployment on edge devices.

In mLSTM, the memory component is non-parametric but involves computationally intensive operations through a $d \times d$ matrix structure for both storage and updating. This approach sacrifices computational simplicity to achieve larger memory capacity. Nonetheless, because these matrix operations are amenable to parallelization on GPUs, their impact on overall runtime is minimal.

mLSTM offers a degree of parallelism comparable to frameworks like FlashAttention~\citep{Dao:22,Dao:23} or GLA~\citep{Yang:23arxiv}, whereas sLSTM does not parallelize efficiently due to its reliance on memory mixing (i.e., hidden-to-hidden connections). To address this, we have engineered an efficient CUDA-based implementation with register-level GPU memory optimization; as a result, sLSTM runs at a speed generally no more than twice as slow as mLSTM.

\section{Related Work}
\label{sec:related}

\paragraph{Linear Attention.} Numerous approaches have been developed to address the Transformer’s quadratic scaling with context length by achieving linear computational complexity. The Synthesizer model generates synthetic attention scores independently of token--token relationships~\citep{Tay:20arxiv}. Linformer employs a low-rank approximation for self-attention, facilitating a linear complexity scheme~\citep{Wang:20arxiv}. Linear Transformer reformulates the attention mechanism to operate in linear time with respect to sequence length~\citep{Katharopoulos:20}. Performer approximates the attention softmax function using a positive orthogonal random features method~\citep{Choromanski:21}. In the Structured Global Convolution (SGConv)~\citep{Li:22} as well as in Hyena Hierarchy~\citep{Poli:23}, attention modules are replaced by efficient long-range convolutional operations.

\paragraph{State Space Models.} State Space Models (SSMs) have recently surged in popularity due to their linear dependence on context length and competitive performance when compared with Transformer architectures. The initial breakthrough came with the introduction of the Structured State Space sequence model (S4)~\citep{Gu:21}. Subsequently, various extensions and alternatives emerged, such as the Diagonal State Space (DSS) model~\citep{Gupta:22}, the Gated State Space (GSS) models~\citep{Mehta:22}, the S5 model~\citep{Smith:22}, Bidirectional Gated SSM (BiGS)~\citep{Wang:22}, H3 model~\citep{Fu:23}, and Mamba~\citep{Gu:24arxiv}.

\paragraph{Recurrent Neural Networks.} Recurrent Neural Networks (RNNs) have been proposed as alternatives to Transformer and attention-based models, owing to their linear scalability with respect to sequence length. Language modeling tasks have benefited from RNNs employing Deep Linear Recurrent Units (LRUs)~\citep{Orvieto:23, De:24arxiv}, while methods such as Hierarchically Gated Linear RNN (HGRN)~\citep{Qin:23} and its successor HGRN2~\citep{Qin:24arxiv} have also demonstrated strong results. Among large language models, RWKV~\citep{Peng:23arxivshort,Peng:24arxivshort} is a prominent RNN-based approach that achieves performance comparable to that of Transformers.

\paragraph{Gating.}
The concept of gating is fundamental to the design of LSTM and has been re-examined and expanded upon in many modern models. A variety of recent architectures have incorporated gating mechanisms, including HGRN~\citep{Qin:23}, HGRN2~\citep{Qin:24arxiv}, Gated Linear Attention (GLA)~\citep{Yang:23arxiv}, Gated State Space (GSS) models~\citep{Mehta:22}, Bidirectional Gated SSM (BiGS)~\citep{Wang:22}, Moving Average Equipped Gated Attention (MEGA)~\citep{Ma:22}, RWKV~\citep{Peng:23arxivshort}, and Mamba~\citep{Gu:24arxiv}.

\paragraph{Covariance Update Rule.}
In an effort to improve memory efficiency, we have introduced a covariance-based matrix update in the mLSTM cell. This strategy of leveraging a covariance update can also be found in Fast Weight Programmers~\citep{Schmidhuber:92ncfastweights,Schlag:21}, RWKV-5 and RWKV-6~\citep{Peng:24arxivshort}, Retention~\citep{Sun:23arxiv}, Linear Transformer~\citep{Katharopoulos:20}, and HGRN2~\citep{Qin:24arxiv}.

\paragraph{Most Related.} The architectures most analogous to xLSTM include Retention~\citep{Sun:23arxiv}, RWKV~\citep{Peng:23arxivshort,Peng:24arxivshort}, and HGRN2~\citep{Qin:24arxiv}. These frameworks employ matrix memory and/or gating mechanisms as central features. Distinct from the recently introduced sLSTM, these prior models lack provisions for memory mixing. Memory mixing is crucial for resolving state tracking tasks, which renders LSTM models more powerful in terms of expressiveness compared to State Space Models (SSMs) and Transformers~\citep{Merrill:24,Deletang:23}. State tracking plays a vital role in applications such as code evaluation and monitoring entities throughout extensive narratives.

\paragraph{Residually Stacking Architectures.} In line with the practices common to modern, large-scale deep learning models, xLSTM architectures employ a residual stacking approach using foundational building blocks~\citep{Srivastava:15,He:16}.

This strategy paved the way for the development of deep convolutional networks~\citep{He:16} and later for Transformers~\citep{Vaswani:17}. The adoption of Transformer architectures has been instrumental in powering Large Language Models (LLMs) such as GPT-3~\citep{Brown:20short}, ChatGPT~\citep{Schulman:22short}, GPT-4~\citep{Achiam:23gpt4short}, Megatron-LM~\citep{Shoeybi:19}, Gopher~\citep{Rae:21short}, ERNIE 3.0 Titan~\citep{Wang:21arxivshort}, GLaM ~\citep{Du:21glamshort}, Chinese M6~\citep{Lin:21}, the multilingual AlexaTM 20B~\citep{Soltan:22}, OPT~\citep{Zhang:22opt}, Chinchilla~\citep{Hoffmann:22short}, BLOOM~\citep{Scao:22arxivshort}, GLM-130B~\citep{Zeng:22arxivshort}, LaMDA~\citep{Thoppilan:22short}, PaLM~\citep{Chowdhery:22short}, Llama~\citep{Touvron:23arxiv}, and Gemini~\citep{GeminiTeam:23arxiv,Reid:24arxiv}.

\section{Experiments}
\label{sec:Exp}

This section presents an empirical evaluation of xLSTM, emphasizing its performance relative to established language modeling techniques. Section~\ref{sec:ExpSyntethic} explores xLSTM's distinctive attributes through synthetic benchmarks. 
In Section~\ref{sec:ExpComparison}, we examine validation set perplexity across various contemporary language modeling strategies, each trained on 15B tokens sourced from SlimPajama~\citep{Soboleva:23}, and perform an ablation analysis specifically for xLSTM on the same dataset. Furthermore, we investigate the scaling dynamics of all considered models following the approach described by \citet{Kaplan:20} and \citet{Brown:20short}.
Section~\ref{sec:ExpLanguage} extends the analysis by conducting comprehensive experiments for language modeling, wherein xLSTM and the leading methods identified in Section~\ref{sec:ExpComparison} are trained on a substantially larger dataset of 300B tokens from SlimPajama~\citep{Soboleva:23}. The evaluation is carried out in four stages: firstly, we determine the capacity of the models to generalize to extended contexts; secondly, their validation perplexity and effectiveness on downstream tasks~\citep{Sutawika:24} are measured; thirdly, we test all models across 571 text domains of the PALOMA language benchmark dataset~\citep{Magnusson:23arxivshort}; lastly, we reassess the scaling behavior, now leveraging a dataset that is twenty times larger.

\label{sec:xlstm_blocks_notation}
Throughout the experiments, we denote xLSTM[$a$:$b$] as the configuration where the ratio $a/b$ indicates the number of mLSTM-based to sLSTM-based xLSTM blocks. As an illustration, xLSTM[7:1] corresponds to a total of eight blocks, where seven are implemented with mLSTM and one with sLSTM. In the standard setting with 48 total blocks, this configuration results in 42 blocks using mLSTM and 6 with sLSTM. Additionally, for all experimental settings, mLSTM employs a pre up-projection block, while sLSTM uses a post up-projection block.

\subsection{Synthetic Tasks and Long Range Arena}
\label{sec:ExpSyntethic}
We begin by evaluating the capabilities of xLSTM's novel exponential gating mechanism combined with memory mixing on a range of formal language benchmarks~\citep{Deletang:23}. Following this, we investigate how effectively the newly introduced matrix memory in xLSTM handles the Multi-Query Associative Recall task~\citep{Arora:23arxiv}. Lastly, we analyze xLSTM's ability to process lengthy sequences within the Long Range Arena suite~\citep{Tay:21}.

\paragraph{Evaluation of xLSTM's Exponential Gating and Memory Mixing Capability.} We investigate the effectiveness of xLSTM's exponential gating mechanism, which incorporates memory mixing, hypothesized to facilitate state tracking tasks~\citep{Merrill:24, Merrill:23}. The formal language benchmarks described in \citet{Deletang:23} are adapted and expanded to support training across variable input lengths, thereby assessing the model's capacity for length generalization. For comprehensive task definitions and supplementary findings, refer to Appendix~\ref{sec:appExpSynthetic-formalLang}. 

xLSTM's performance is juxtaposed with a range of architectures, including Transformers, State Space Models, and Recurrent Neural Networks. Evaluation is carried out by measuring accuracy on task-specific target tokens, with the scores normalized between 0 (random guessing) and 1 (perfect accuracy). The comparative analysis focuses on 2-block variants for each approach, such as xLSTM[0:1] (pure sLSTM), xLSTM[1:0] (pure mLSTM), xLSTM[1:1], Llama, Mamba, RWKV, Retention, Hyena, LSTM, as well as LSTM embedded within Transformer layers (LSTM (Block)). The corresponding outcomes are depicted in Figure~\ref{fig:formal-main}.

Models lacking memory mixing or explicit state tracking—such as standard Transformers or State Space Models—are unable to resolve tasks like regular grammars, exemplified by the parity problem. This observation corroborates previous conclusions that RNNs exhibit fundamentally greater expressivity in such settings compared to Transformers and State Space models~\citep{Merrill:24, Merrill:23, Deletang:23}.

\paragraph{Test of xLSTM's Memory Capacities on Associative Recall Tasks.} This study investigates the matrix memory mechanism introduced in xLSTM by evaluating its ability to store and retrieve information using the Multi-Query Associative Recall task~\citep{Arora:23arxiv}. In this benchmark, sequences are composed of randomly assigned key-value pairs drawn from an extensive vocabulary; these associations must be retained for subsequent access. To further challenge the models beyond the baseline task, the number of key-value pairs is scaled up to 256 and context lengths are extended to 2048 tokens. This setup facilitates a comprehensive examination of memory capacity across various model designs. The evaluation contrasts 2-block versions of Llama, Mamba, RWKV-5, RWKV-6, xLSTM[1:1], and xLSTM[1:0], and measures their retrieval accuracy. Given the exponential scaling of memory capacity with respect to the coding dimension in Transformer-based architectures (e.g., Llama)~\citep{Ramsauer:21}, these serve as the primary benchmark for comparison. Performance metrics are illustrated in Figure~\ref{fig:mqar-main}. Notably, xLSTM[1:1] surpasses all non-Transformer competitors, performing robustly even in smaller model regimes. Interestingly, the inclusion of the sLSTM block does not impede memory retention; on the contrary, it enhances memory utilization, as demonstrated in the setting with 256 key-value pairs—the most taxing scenario considered. Further supporting results, including extrapolation studies, are provided in Appendix~\ref{sec:appExpSynthetic-mqar}, revealing that xLSTM’s improved memory system enables successful generalization to longer contexts beyond those encountered during model training.

\paragraph{Test of xLSTM's Long Context Capabilities on Long Range Arena.} In order to evaluate the effectiveness of xLSTM with extended sequences and broad contextual information, we conduct a comparative analysis of several approaches using Long Range Arena~\citep{Tay:21}. Our results show that xLSTM reliably achieves high performance across each task, indicating that the architecture excels at managing diverse long-context challenges with notable efficiency.

For more details, see Appendix~\ref{sec:appExpSynthetic-lra}.

\subsection{Method Comparison and Ablation Study}
\label{sec:ExpComparison}

In order to explore the central inquiry of this work—namely, the capabilities of our newly proposed LSTM variants in large-scale language modeling—we conduct training experiments with xLSTMs, Transformers, State Space Models, and alternative approaches on a dataset of 15B tokens from SlimPajama, utilizing an identical auto-regressive training scheme. We evaluate model performance on the validation dataset, and further analyze the xLSTMs through comprehensive ablation experiments.

\paragraph{Benchmarking xLSTM Against Alternative Approaches.} To establish a baseline for comparison, we trained various models using a dataset of 15B tokens from SlimPajama~\citep{Soboleva:23}. Model performance was quantified using perplexity scores on a designated validation set. The evaluation encompassed multiple architectures: our proposed xLSTM, GPT-3 (Transformer)~\citep{Brown:20short}, Llama (Transformer)~\citep{Touvron:23arxiv}, H3 (SSM)~\citep{Fu:23}, Mamba (SSM)~\citep{Gu:24arxiv}, RWKV-4, RWKV-5, and RWKV-6 (all RNNs)~\citep{Peng:23arxivshort, Peng:24arxivshort}, GLA (linear Transformer)~\citep{Yang:23arxiv}, HGRN and HGRN2 (RNNs)~\citep{Qin:23, Qin:24arxiv}, as well as RetNet (linear Transformer)~\citep{Sun:23arxiv}, Hyena (linear Transformer)~\citep{Poli:23}, xLSTM[1:0], and xLSTM[7:1] (see Section~\ref{sec:xlstm_blocks_notation}). 

All models were trained using mixed precision. Notably, RWKV-5, RWKV-6, GLA, and HGRN2 employed a mixed precision approach that did not rely on PyTorch's automated mixed precision tools (refer to Appendix Section~\ref{sec:appGenTrainProc} for details). We organized the methods into three groups: (a) Transformer-based models, (b) State Space Models (SSMs), and (c) Recurrent Neural Networks (RNNs) and linear Transformers. Here, linear Transformers leverage linear alternatives to the standard attention mechanism. Each model was constructed to align with a GPT-3 configuration comprising 350M parameters—embedding dimension of 1024 and 24 residual blocks. Only GPT-3 utilizes tied weights for token and output embeddings, resulting in a reduced parameter count. 

As indicated in Table~\ref{tab:spaj15b_model_results}, xLSTM achieves superior validation perplexity compared to all competing methods. Further experimental information is provided in Appendix~\ref{sec:appExpComparison}. Additionally, Figure~\ref{fig:temp_sclaw15B} presents scaling trends for these experiments, suggesting that xLSTM maintains an advantage as model size increases.

\paragraph{Ablation Studies.}
As presented in Table~\ref{tab:spaj15b_model_results} and Figure~\ref{fig:temp_sclaw15B}, xLSTM attains remarkable performance in language modeling when trained on 15B tokens sourced from SlimPajama. To dissect the contributions of each modification transitioning from the conventional LSTM to xLSTM, we incrementally transform the standard LSTM architecture into that of xLSTM. The initial step involves embedding LSTM layers within pre-LayerNorm residual networks. The next stage incorporates an up-projection block subsequent to these layers. Ultimately, we introduce exponential gating mechanisms and matrix memory to complete the xLSTM framework.

The results are shown in Table~\ref{tab:ablstudies} (top). 

Ablation studies indicate that the observed performance gains result from the combined contributions of exponential gating and matrix memory. Moreover, given the central role of gating in RNNs and State Space Models, we investigate the impact of various gating strategies. As shown in Table~\ref{tab:ablstudies} (bottom), the results demonstrate that configuring each gate to be trainable and input-dependent yields progressively better outcomes. Further detailed analysis concerning the separate backbone elements is provided in Appendix~\ref{sec:appAblDetails}.

\subsection{xLSTM as Large Language Model}
\label{sec:ExpLanguage}

This work concludes with comprehensive language modeling experiments at scale, aiming to evaluate xLSTM as a candidate for LLM applications. To this end, we significantly increase the volume of training data, utilizing 300B tokens from SlimPajama. Notably, this token count aligns with the datasets employed by Mamba~\citep{Gu:24arxiv} and Griffin~\citep{De:24arxiv}.

For benchmarking, xLSTM is contrasted against RWKV-4, Llama, and Mamba—each selected to represent different method classes outlined in Section~\ref{sec:ExpComparison}. RWKV-4 is chosen as the RNN baseline, since suitable training precision settings for RWKV-5, RWKV-6, and HGRN2 were established only after initiating the 300B token training runs (further details in Appendix~\ref{sec:appGenTrainProc}).

A range of model configurations (125M, 350M, 760M, 1.3B parameters) are trained and systematically evaluated for sequence length extrapolation and validation performance. Additional evaluations include downstream task proficiency, language modeling results across 571 text domains from the PALOMA benchmark, and an analysis of their scaling law characteristics.

\paragraph{Sequence Length Extrapolation.}
\label{sec:spaj300B_ctx_extrapolation}
To begin, we investigate the capability of large-scale models (1.3B parameters)---specifically xLSTM, RWKV-4, Llama, and Mamba---to generalize to sequences longer than those seen during training. Each model is trained with a maximum context length of 2048 tokens and subsequently assessed at context lengths extending to 16384 tokens. Performance outcomes are depicted in Figure~\ref{fig:spaj300B_seq_extrapolation}. Unlike the alternative architectures, xLSTM consistently achieves low perplexity values even as the context length increases.

\paragraph{Validation Perplexity and Downstream Tasks.}
\label{sec:spaj_downstream_eval}
Next, across all considered model scales, we benchmark xLSTM, RWKV-4, Llama, and Mamba on the SlimPajama validation dataset for next token prediction as well as on a suite of downstream evaluations designed to test models’ ability in common sense reasoning.

The third column of Table~\ref{tab:lmeval} reports validation set perplexity scores for each evaluated method. Across all model sizes, both xLSTM[1:0] and xLSTM[7:1] achieve the lowest perplexity on the validation set. The remaining columns in Table~\ref{tab:lmeval} summarize the results obtained on various downstream tasks. In almost all cases and for every model size considered, xLSTM outperforms the other approaches; Mamba exhibits superior performance over xLSTM only in selected instances of the ARC task. Further discussion is provided in Appendix~\ref{sec:appExpLanguage}.

\paragraph{Performance on PALOMA Language Tasks.} Next, across various model scales, we evaluate xLSTM, RWKV-4, Llama, and Mamba on the PALOMA language tasks~\citep{Magnusson:23arxivshort} using next token prediction. We assess model performance based on perplexity scores calculated over 571 distinct text domains, which encompass sources like nytimes.com and the Reddit community r/depression.
Token prediction perplexity results are presented in Table~\ref{tab:ppleval}, with language modeling domains listed in the initial seven columns, and specific domain benchmarks grouped in the final five. For these tasks, xLSTM[1:0] achieves superior performance compared to xLSTM[7:1].

Specifically, xLSTM[1:0] demonstrates a lower perplexity than Mamba in 568 out of 571 text domains (99.5\%), outperforms Llama in 486 out of 571 domains (85.1\%), and shows improved perplexity compared to RWKV-4 in 570 out of 571 domains (99.8\%). Additional details are provided in Appendix~\ref{sec:appExpLanguage}.

\paragraph{Scaling Laws.} Next, we investigate the power-law scaling phenomenon, which enables predicting model performance as the size increases~\citep{Kaplan:20,Brown:20short}. The scaling trends are illustrated in Figure~\ref{fig:temp_sclaw300B}. Although all evaluated architectures display comparable scaling characteristics, they differ in their baseline performance levels. RWKV-4 exhibits the lowest performance, with Llama and Mamba outperforming it. Notably, xLSTM surpasses Mamba, with the difference in performance between xLSTM and Mamba being roughly equivalent to that between Mamba and Llama. These results suggest that as model sizes grow, xLSTM maintains a superior scaling advantage relative to both Transformer and State-Space architectures.

\textbf{Generation Times and Maximal Throughput.} We assess the duration required for text generation in Figure~\ref{fig:inference_time_generation} and report the peak throughput results in Figure~\ref{fig:inference_time_decoding_throughput} (left) using various xLSTM models, each at the 1.3B parameter count. These are evaluated alongside comparably sized Mamba, Llama, and RWKV implementations from HuggingFace; for Llama, a static key-value cache is utilized. At the point of conducting these benchmarks, it was not possible to compile the full cache for the Transformer Model or to use \texttt{torch.compile} to compile Mamba. For the generation benchmarks, every model is evaluated with a batch size of 1 and a pre-fill value of 16, creating the most advantageous conditions for the Transformer architecture. 

As illustrated in Figure~\ref{fig:inference_time_generation}, both xLSTM and the recurrent Mamba and RWKV-4 models display linear growth in generation time, in contrast to the quadratic increase seen in the Llama model. Throughput analysis considers a variety of batch sizes and pre-fill values specifically for Llama. Finally, Figure~\ref{fig:inference_time_decoding_throughput} (right) highlights that xLSTM supports considerably larger batch sizes compared to Llama, owing to xLSTM's constant memory usage, resulting in superior maximal throughput.

\section{Limitations}

(i) Unlike mLSTM, the structure of sLSTM involves memory mixing that restricts parallelizable computations and thus precludes efficient parallel processing. Despite this, we implemented an optimized CUDA kernel for sLSTM that currently operates at less than twice the runtime of our parallel mLSTM kernel. (ii) At present, our mLSTM CUDA kernels lack specific optimizations, making their execution approximately four times slower compared to FlashAttention and the scan method used in Mamba. Adopting techniques similar to those in FlashAttention would likely yield much faster CUDA kernels. (iii) The matrix memory utilized by mLSTM incurs substantial computational demands because it handles $d \times d$ matrix operations. However, both updating and retrieving memory are parameter-free and amenable to parallelization through conventional matrix computations, so the actual wall-clock overhead from complex memory access is minimal. (iv) Proper initialization of forget gates is critical to effective operation. (v) Since matrix memory does not scale with sequence length, longer sequences may saturate memory capacity when extended contexts are considered. Nevertheless, this issue has not surfaced for contexts up to 16k tokens, as detailed in Section~\ref{sec:spaj300B_ctx_extrapolation}. (vi) Because large language experiments impose significant computational costs, neither the architecture nor the hyperparameters—especially for larger xLSTM configurations—have undergone thorough optimization. Comprehensive tuning is expected to be necessary for xLSTM to realize its maximal performance.

\section{Conclusion}
Our initial inquiry asked: How effective is language modeling with LSTM architectures scaled to billions of parameters? Based on our findings, the answer is: “They perform at least as well as leading methods such as Transformers or State Space Models.” We extended the conventional LSTM by introducing exponential gating, memory mixing, and a novel memory configuration—collectively termed xLSTM. xLSTM demonstrates superior performance in language modeling when benchmarked against top-tier approaches like Transformers and State Space Models. Scaling behavior suggests that, as xLSTM models grow in size, they are likely to become strong contenders against large language models built using Transformers. Beyond language modeling, xLSTM presents promising prospects for applications in domains such as Reinforcement Learning, Time Series Analysis, and modeling complex physical phenomena.

\end{document}